% carried verbatim from the v2 tree (hardware.tex); NOT part of this session's deliverable.
\begin{figure*}[tp]\centering% self-contained native pgfplots fragment (fig:noiserec)
\definecolor{nrH6}{HTML}{6A4C93}
\definecolor{nrNH3}{HTML}{E8880C}
\definecolor{nrN2}{HTML}{D7263D}
\definecolor{nrCO}{HTML}{12876F}
\definecolor{nrC2}{HTML}{2E6F95}
\definecolor{nrOK}{HTML}{0E7C7B}\definecolor{nrBad}{HTML}{5E4B8B}\definecolor{nrNv}{HTML}{D1495B}\definecolor{nrChem}{HTML}{2E7D32}
\begin{tikzpicture}
\begin{groupplot}[group style={group size=2 by 1, group name=nrg, horizontal sep=1.45cm}, height=9.5cm, paperaxis,
  title style={at={(0,1)},anchor=south west,font=\bfseries,yshift=1pt}]
\nextgroupplot[width=0.505\linewidth, ymode=log, xmin=-1, xmax=31, ymin=3e-3, ymax=80,
  xlabel={per-qubit bit-flip rate $\varepsilon$ (\%)}, ylabel={energy error vs FCI (mHa)}, title={(a)},
  legend style={at={(0.5,0.03)},anchor=south,font=\normalsize,draw=black!20},legend columns=3,legend cell align=left]
\draw[nrChem,dashed,line width=0.6pt] (axis cs:-1,1.6) -- (axis cs:31,1.6);
% chemical-accuracy label: emitted after the groupplot, see below.
\addplot[name path=H6lo,draw=none,forget plot] coordinates {(0,3.2914) (2,3.0066) (5,3.6550) (10,3.7659) (15,5.3712) (20,8.9720) (30,15.6267)};
\addplot[name path=H6hi,draw=none,forget plot] coordinates {(0,4.5041) (2,4.1385) (5,4.4412) (10,6.5199) (15,9.5780) (20,16.7961) (30,29.8587)};
\addplot[nrH6,opacity=0.15,forget plot] fill between[of=H6lo and H6hi];
\addplot[name path=NH3lo,draw=none,forget plot] coordinates {(0,2.1436) (2,1.4495) (5,1.0457) (10,0.7919) (15,1.3563) (20,1.8046) (30,5.7326)};
\addplot[name path=NH3hi,draw=none,forget plot] coordinates {(0,2.9078) (2,1.9853) (5,1.6281) (10,1.6847) (15,2.0580) (20,2.8561) (30,20.1868)};
\addplot[nrNH3,opacity=0.15,forget plot] fill between[of=NH3lo and NH3hi];
\addplot[name path=N2lo,draw=none,forget plot] coordinates {(0,0.1729) (2,0.1550) (5,0.3169) (10,0.2965) (15,0.6158) (20,1.4642) (30,8.4636)};
\addplot[name path=N2hi,draw=none,forget plot] coordinates {(0,0.6006) (2,0.6093) (5,0.6846) (10,1.3158) (15,2.3324) (20,5.5026) (30,17.8734)};
\addplot[nrN2,opacity=0.15,forget plot] fill between[of=N2lo and N2hi];
\addplot[name path=COlo,draw=none,forget plot] coordinates {(0,1.5510) (2,1.5669) (5,1.0491) (10,0.7344) (15,0.5781) (20,0.4698) (30,0.3702)};
\addplot[name path=COhi,draw=none,forget plot] coordinates {(0,1.9155) (2,2.0418) (5,1.4264) (10,1.3026) (15,1.0358) (20,0.7942) (30,1.0249)};
\addplot[nrCO,opacity=0.15,forget plot] fill between[of=COlo and COhi];
\addplot[name path=C2lo,draw=none,forget plot] coordinates {(0,0.1703) (2,0.1351) (5,0.0827) (10,0.0677) (15,0.0479) (20,0.0385) (30,0.0118)};
\addplot[name path=C2hi,draw=none,forget plot] coordinates {(0,0.1967) (2,0.1807) (5,0.1240) (10,0.0817) (15,0.0730) (20,0.0627) (30,0.0433)};
\addplot[nrC2,opacity=0.15,forget plot] fill between[of=C2lo and C2hi];
\addplot[nrH6,line width=1.4pt,mark=*,mark size=1.3pt,mark options={fill=nrH6,draw=white,line width=0.4pt}] coordinates {(0,3.8978) (2,3.5726) (5,4.0481) (10,5.1429) (15,7.4746) (20,12.8840) (30,22.7427)}; \addlegendentry{H$_6$}
\addplot[nrH6,line width=0.9pt,densely dashed,forget plot] coordinates {(0,3.8978) (2,4.6251) (5,5.6347) (10,7.8697) (15,12.9901) (20,23.8560) (30,50.3837)};
\addplot[nrNH3,line width=1.4pt,mark=*,mark size=1.3pt,mark options={fill=nrNH3,draw=white,line width=0.4pt}] coordinates {(0,2.5257) (2,1.7174) (5,1.3369) (10,1.2383) (15,1.7072) (20,2.3304) (30,12.7391)}; \addlegendentry{NH$_3$}
\addplot[nrNH3,line width=0.9pt,densely dashed,forget plot] coordinates {(0,2.5257) (2,2.9622) (5,2.7617) (10,2.9227) (15,3.7294) (20,5.5505) (30,31.5152)};
\addplot[nrN2,line width=1.4pt,mark=*,mark size=1.3pt,mark options={fill=nrN2,draw=white,line width=0.4pt}] coordinates {(0,0.3843) (2,0.3445) (5,0.5007) (10,0.6589) (15,1.3685) (20,3.2538) (30,13.1685)}; \addlegendentry{N$_2$}
\addplot[nrN2,line width=0.9pt,densely dashed,forget plot] coordinates {(0,0.3843) (2,0.4335) (5,0.6929) (10,1.1118) (15,2.8006) (20,6.3706) (30,31.8749)};
\addplot[nrCO,line width=1.4pt,mark=*,mark size=1.3pt,mark options={fill=nrCO,draw=white,line width=0.4pt}] coordinates {(0,1.7332) (2,1.8043) (5,1.2377) (10,1.0185) (15,0.8069) (20,0.6320) (30,0.6975)}; \addlegendentry{CO}
\addplot[nrCO,line width=0.9pt,densely dashed,forget plot] coordinates {(0,1.7332) (2,1.9851) (5,1.9371) (10,1.4901) (15,1.8518) (20,2.5550) (30,2.8339)};
\addplot[nrC2,line width=1.4pt,mark=*,mark size=1.3pt,mark options={fill=nrC2,draw=white,line width=0.4pt}] coordinates {(0,0.1835) (2,0.1579) (5,0.1034) (10,0.0747) (15,0.0604) (20,0.0506) (30,0.0261)}; \addlegendentry{C$_2$}
\addplot[nrC2,line width=0.9pt,densely dashed,forget plot] coordinates {(0,0.1835) (2,0.1635) (5,0.1143) (10,0.0811) (15,0.0746) (20,0.0647) (30,0.0657)};
\nextgroupplot[width=0.425\linewidth, xmode=log, xmin=2.2e-3, xmax=14, enlarge y limits=0.03,
  ytick={0,1,2,3,4,5,6,7,8,9,10,11,12,13,14,15,16,17,18}, yticklabels={BeH$_2$,H$_2$,LiH,HF,F$_2$,H$_2$O,H$_4$,OH,LiF,BH,N$_2$,BeH,C$_2$,NO,O$_2$,CN,CO,NH$_3$,H$_6$}, y tick label style={font=\normalsize},
  xtick={0.01,0.1,1,10}, xticklabels={$0.01$,$0.1$,$1$,$10$}, xminorticks=false,
  xlabel={energy error at $\varepsilon=2\%$ (mHa)}, title={(b)},
  legend style={at={(0.03,0.03)},anchor=south west,font=\normalsize,draw=black!20},legend cell align=left]
\fill[nrChem,fill opacity=0.06] (axis cs:2.2e-3,-0.6) rectangle (axis cs:1.6,18.6);
\draw[nrChem,dashed,line width=0.7pt] (axis cs:1.6,-0.6) -- (axis cs:1.6,18.6);
\node[nrChem,font=\tiny,anchor=south,rotate=90] at (axis cs:1.6,9.5) {};
\draw[black!32,line width=0.8pt] (axis cs:0.0022,0) -- (axis cs:0.0022,0); \draw[black!32,line width=0.8pt] (axis cs:0.0022,1) -- (axis cs:0.0022,1); \draw[black!32,line width=0.8pt] (axis cs:0.0022,2) -- (axis cs:0.0022,2); \draw[black!32,line width=0.8pt] (axis cs:0.0022,3) -- (axis cs:0.0022,3); \draw[black!32,line width=0.8pt] (axis cs:0.0022,4) -- (axis cs:0.0022,4); \draw[black!32,line width=0.8pt] (axis cs:0.0327,5) -- (axis cs:0.0022,5); \draw[black!32,line width=0.8pt] (axis cs:0.3178,6) -- (axis cs:0.0022,6); \draw[black!32,line width=0.8pt] (axis cs:0.0194,7) -- (axis cs:0.0194,7); \draw[black!32,line width=0.8pt] (axis cs:0.0254,8) -- (axis cs:0.0254,8); \draw[black!32,line width=0.8pt] (axis cs:0.0845,9) -- (axis cs:0.0605,9); \draw[black!32,line width=0.8pt] (axis cs:0.0688,10) -- (axis cs:0.0636,10); \draw[black!32,line width=0.8pt] (axis cs:0.1367,11) -- (axis cs:0.1367,11); \draw[black!32,line width=0.8pt] (axis cs:0.1711,12) -- (axis cs:0.1615,12); \draw[black!32,line width=0.8pt] (axis cs:0.3275,13) -- (axis cs:0.1942,13); \draw[black!32,line width=0.8pt] (axis cs:0.4696,14) -- (axis cs:0.3363,14); \draw[black!32,line width=0.8pt] (axis cs:0.7815,15) -- (axis cs:0.7261,15); \draw[black!32,line width=0.8pt] (axis cs:1.3923,16) -- (axis cs:1.2450,16); \draw[black!32,line width=0.8pt] (axis cs:2.6491,17) -- (axis cs:1.5063,17); \draw[black!32,line width=0.8pt] (axis cs:4.3814,18) -- (axis cs:3.5606,18);
\addplot[only marks,mark=o,mark size=1.9pt,nrNv,line width=0.8pt] coordinates {(0.0022,0) (0.0022,1) (0.0022,2) (0.0022,3) (0.0022,4) (0.0327,5) (0.3178,6) (0.0194,7) (0.0254,8) (0.0845,9) (0.0688,10) (0.1367,11) (0.1711,12) (0.3275,13) (0.4696,14) (0.7815,15) (1.3923,16) (2.6491,17) (4.3814,18)}; \addlegendentry{naive}
\addplot[only marks,mark=*,mark size=2.1pt,nrOK] coordinates {(0.0022,0) (0.0022,1) (0.0022,2) (0.0022,3) (0.0022,4) (0.0022,5) (0.0022,6) (0.0194,7) (0.0254,8) (0.0605,9) (0.0636,10) (0.1367,11) (0.1615,12) (0.1942,13) (0.3363,14) (0.7261,15) (1.2450,16) (1.5063,17) (3.5606,18)}; \addlegendentry{S-CoRe}
\node[black,font=\normalsize,anchor=north east] at (axis cs:1.4720000000000002,18.3) {$1.6$ mHa};
\end{groupplot}
\node[nrChem,font=\footnotesize,anchor=south east,inner sep=1pt]
      at (nrg c1r1.north east) {chemical accuracy ($1.6$\,mHa)};
\end{tikzpicture}

\caption{\textbf{SQD ground-state energy under a controlled bit-flip model across the molecular suite.}
\textbf{(a)}~Energy error vs the per-qubit bit-flip rate $\varepsilon$ for five representative molecules
($\mathrm{H_6}$, $\mathrm{NH_3}$, $\mathrm{N_2}$, CO, $\mathrm{C_2}$; seed-averaged over $8$ seeds at $4{,}000$ shots each, bands $=\pm1$ s.d.\ across seeds; log scale). Solid $=$ S-CoRe recovery, dashed $=$ naive
post-selection. The advantage of recovery over naive post-selection, both arms consuming the same corrupted
bitstrings at the same rate with the same seeds, grows with noise. For $\mathrm{CO}$ and
$\mathrm{C_2}$ the error at finite $\varepsilon$ also falls below its own $\varepsilon=0$ value;
that second comparison is not used here (Sec.~\ref{sec:scope:hw}): it is decided by the size of
the recovered subspace, which this sweep does not record. \textbf{(b)}~The full nineteen-molecule suite at a realistic
$\varepsilon=2\%$ (log scale, sorted): for each molecule the naive post-selection error (open circle) is
linked to the S-CoRe value (filled), so the leftward pull is the recovery gain; the shaded band is
chemical accuracy ($<1.6$~mHa). Eighteen of nineteen land inside it; $\mathrm{H_6}$, the most
multireference case, is the residual challenge. Panels (a) and (b) are independent seed ensembles
(panel (a): $8$ seeds; panel (b): $5$ seeds, $4000$ shots), so a given molecule's $\varepsilon=2\%$ value
can differ by up to a factor of a few between them (for $\mathrm{N_2}$, $0.34$ vs $0.06$~mHa; both far below chemical accuracy).}\label{fig:noiserec}\end{figure*}
